\section{Extremal primitives and boundary readout on boxes}

We now turn to the main equality-and-stability theorem chain. The sharp lower
bound in \eqref{eq:intro-extremal-problem} is governed by the crystalline
perimeter
\[
P_1(E):=\int_{\partial E}\abs{\nu_E}_1\,\dd\Haus^{n-1},
\qquad
\abs{\nu_E}_1:=\sum_{i=1}^n \abs{(\nu_E)_i}.
\]
In the preliminary general-domain statements in this section the ambient
dimension is the displayed dimension $n$. In the rectangular-box theorems the
dimension is the active top-form dimension $k$ of
$R=\prod_{j=1}^k[0,L_j]$; the same notation $P_1$ and $C_\infty$ is used with
that displayed dimension.
For background on crystalline isoperimetry, the Wulff theorem, and calibrable
sets we refer to
\cite{Taylor1978,Fonseca1991,FonsecaMuller1991,DacorognaPfister1992,
FigalliZhang2022,DeMason2024} and
\cite{Strang1983,Nozawa1990,Grieser2006,Strang2010,
CasellesFaccioloMeinhardt2009,CasellesChambolleMollNovaga2008,
Leonardi2015,Parini2011}.

\begin{definition}[The coefficient-$L^\infty$ divergence constant]
\label{def:C-infty}
Let $\Omega\subset\RR^n$ be a bounded Lipschitz domain and set
$\mathrm{vol}:=\dd x_1\wedge\cdots\wedge\dd x_n$. Define
\[
C_\infty(\Omega):=
\inf\left\{
\norm{V}_{L^\infty(\Omega;\ell^\infty)}:
\begin{array}{l}
\overline\Omega\subset U\subset\RR^n\ \text{with }U\text{ open},\\
V\in C^\infty(U;\RR^n),\\
\Div V=1\ \text{on }\Omega
\end{array}
\right\}.
\]
Equivalently, $C_\infty(\Omega)$ is the infimum of
$\norm{\eta}_{\coeff,\infty}$ over all smooth $(n-1)$-forms $\eta$ defined on a
neighborhood of $\Omega$ and satisfying $\dd\eta=\mathrm{vol}$ on $\Omega$.
The equivalence is through
\begin{equation}
\label{eq:form-vector-identification}
\eta=
\sum_{j=1}^n(-1)^{j-1}V_j\,
\dd x_1\wedge\cdots\wedge\widehat{\dd x_j}\wedge\cdots\wedge\dd x_n,
\end{equation}
for which $\dd\eta=(\Div V)\mathrm{vol}$ and
$\norm{\eta}_{\coeff,\infty}=\norm{V}_{L^\infty(\Omega;\ell^\infty)}$.
\end{definition}

\begin{lemma}[Anisotropic calibration bound]
\label{lem:boundary-calibration}
Let $\Omega\subset\RR^n$ be a bounded Lipschitz domain. If
$V\in C^\infty(U;\RR^n)$, $U\supset\overline\Omega$, and $\Div V=1$ on
$\Omega$, then every Lipschitz subdomain $E\subset\Omega$ satisfies
\begin{equation}
\label{eq:anisotropic-calibration-bound}
\Vol(E)\le \norm{V}_{L^\infty(\Omega;\ell^\infty)}P_1(E).
\end{equation}
Consequently,
\begin{equation}
\label{eq:C-infty-lower-bound}
C_\infty(\Omega)\ge
\sup\left\{\frac{\Vol(E)}{P_1(E)}:\ E\subset\Omega,\ E\ {\rm Lipschitz}\right\}.
\end{equation}
\end{lemma}

\begin{proof}
Let $M:=\norm{V}_{L^\infty(\Omega;\ell^\infty)}$. Since $\Div V=1$ on $E$,
the divergence theorem for Lipschitz domains gives
\[
\Vol(E)=\int_E\Div V\,\dd x
=\int_{\partial E}V\cdot\nu_E\,\dd\Haus^{n-1}.
\]
Pointwise on $\partial E$,
\[
V\cdot\nu_E
\le
\abs{V\cdot\nu_E}
\le
\sum_{j=1}^n\abs{V_j}\abs{(\nu_E)_j}
\le
M\abs{\nu_E}_1 .
\]
Integrating gives \eqref{eq:anisotropic-calibration-bound}. Taking the
infimum over all admissible $V$ and then the supremum over Lipschitz
subdomains $E$ proves \eqref{eq:C-infty-lower-bound}.
\end{proof}

\begin{proposition}[Quantitative boundary saturation on optimal domains]
\label{prop:quantitative-boundary-saturation}
Let $\Omega\subset\RR^n$ be a bounded Lipschitz domain and assume that
\[
C_\infty(\Omega)=\frac{\Vol(\Omega)}{P_1(\Omega)}.
\]
Let $V\in C^\infty(U;\RR^n)$ be admissible for $C_\infty(\Omega)$, meaning
$U\supset\overline\Omega$ is open and $\Div V=1$ on $\Omega$, and set
\[
M:=\norm{V}_{L^\infty(\Omega;\ell^\infty)}.
\]
Then the exact boundary deficit identity
\begin{equation}
\label{eq:boundary-deficit-identity}
\int_{\partial\Omega}
\bigl(M\abs{\nu_\Omega}_1-V\cdot\nu_\Omega\bigr)\,\dd\Haus^{n-1}
=
P_1(\Omega)\bigl(M-C_\infty(\Omega)\bigr)
\end{equation}
holds. In particular, if $M=C_\infty(\Omega)+\varepsilon$, then
\[
\norm{M\abs{\nu_\Omega}_1-V\cdot\nu_\Omega}_{L^1(\partial\Omega)}
=
\varepsilon\,P_1(\Omega).
\]
Every minimizer $V$ for $C_\infty(\Omega)$ therefore satisfies
\[
V\cdot\nu_\Omega
=C_\infty(\Omega)\abs{\nu_\Omega}_1
\quad\text{for }\Haus^{n-1}\text{-a.e. }x\in\partial\Omega .
\]
This conclusion is only almost everywhere; pointwise face traces below require
the additional continuity argument in Theorem~\ref{thm:box-boundary-stability}.
\end{proposition}

\begin{proof}
The integrand in \eqref{eq:boundary-deficit-identity} is nonnegative by
$V\cdot\nu_\Omega\le M\abs{\nu_\Omega}_1$. Hence its $L^1$ norm is its
integral. The divergence theorem and the optimal-domain hypothesis give
\[
\int_{\partial\Omega}V\cdot\nu_\Omega\,\dd\Haus^{n-1}
=\int_\Omega\Div V\,\dd x
=\Vol(\Omega)
=C_\infty(\Omega)P_1(\Omega).
\]
Also
\[
\int_{\partial\Omega}M\abs{\nu_\Omega}_1\,\dd\Haus^{n-1}
=M P_1(\Omega).
\]
Subtracting these two identities proves
\eqref{eq:boundary-deficit-identity}. If $M=C_\infty(\Omega)$, the
nonnegative integrand has integral zero and hence vanishes almost everywhere.

\end{proof}

\begin{lemma}[One-dimensional zero-extension perimeter]
\label{lem:one-dimensional-zero-extension}
Let $A\subset[0,L]$ have finite perimeter, with $L>0$, and let
$\widetilde{\mathbf 1_A}$ be the extension of $\mathbf 1_A$ by zero to
$\RR$. Then
\[
|D\widetilde{\mathbf 1_A}|(\RR)\ge \frac{2}{L}\abs{A}.
\]
This includes the null case $\abs{A}=0$ and the full-measure case
$\abs{A}=L$, where the two jumps are the endpoint jumps of the zero extension.
\end{lemma}

\begin{proof}
We use the standard one-dimensional BV structure theorem for characteristic
functions: after modification on a null set, a measurable subset of an interval
is an at most countable disjoint union of relative open intervals, and the total
variation of its zero extension is the counting measure of the jump endpoints
of those intervals, including jumps at $0$ and $L$ when the zero extension is
viewed on $\RR$ \cite[Sec.~5.10]{EvansGariepy2015}.  In the present
finite-perimeter case this endpoint count is finite whenever it is nonzero.
Thus, if $A$ is null, both sides vanish. If $\abs{A}>0$, at least one interval
component occurs and the zero extension has a transition from $0$ to $1$ and a
transition back from $1$ to $0$, possibly at the endpoints $0$ and $L$. Hence
$|D\widetilde{\mathbf 1_A}|(\RR)\ge2$. Since $\abs{A}\le L$, we obtain
$(2/L)\abs{A}\le2\le |D\widetilde{\mathbf 1_A}|(\RR)$.
\end{proof}

\begin{lemma}[Axis-aligned boxes attain the crystalline ratio]
\label{lem:box-self-cheeger}
Let
\[
R:=\prod_{j=1}^k[0,L_j],
\qquad
m_R:=\frac{1}{2\sum_{j=1}^k L_j^{-1}},
\]
where all $L_j>0$. Then every Lipschitz subdomain $E\subset R$ satisfies
\begin{equation}
\label{eq:box-anisotropic-ineq}
2\left(\sum_{j=1}^k\frac1{L_j}\right)\Vol(E)\le P_1(E).
\end{equation}
Consequently,
\[
\sup\left\{\frac{\Vol(E)}{P_1(E)}:\ E\subset R,\ E\ {\rm Lipschitz}\right\}
=
\frac{\Vol(R)}{P_1(R)}
=m_R,
\]
and
\[
C_\infty(R)=m_R.
\]
Moreover, the affine vector field
\[
V_R(x):=
2m_R\sum_{j=1}^k
\frac{x_j-L_j/2}{L_j}\,\partial_{x_j}
\]
is admissible and satisfies
\[
\norm{V_R}_{L^\infty(R;\ell^\infty)}=m_R.
\]
\end{lemma}

\begin{proof}
Throughout this proof $\mathbf 1_E$ is extended by zero outside $R$, and
$P_1(E)$ denotes the full ambient anisotropic perimeter in $\RR^k$, including
all trace contributions on $\partial R$.
Fix $j$ and write $x=(x_j,x')$ with
$x'\in R'_j:=\prod_{i\ne j}[0,L_i]$. Since $E$ is Lipschitz, it has finite
perimeter, and the BV slicing theorem for sets of finite perimeter
\cite[Sec.~5.10]{EvansGariepy2015} gives, for almost every $x'\in R'_j$, a
finite-perimeter one-dimensional slice
\[
E_{x'}^j:=\{s\in[0,L_j]:(s,x')\in E\}
\]
and the sliced variation identity
\begin{equation}
\label{eq:box-sliced-variation}
\int_{R'_j}P(E_{x'}^j;\RR)\,\dd x'
=|D_{x_j}\mathbf 1_E|(\RR^k)
=\int_{\partial E}\abs{(\nu_E)_j}\,\dd\Haus^{k-1}.
\end{equation}
Here the slices are extended by zero outside $[0,L_j]$, so the relative
end-point jumps at $0$ and $L_j$ are included in $P(E_{x'}^j;\RR)$.
Lemma~\ref{lem:one-dimensional-zero-extension} gives
\[
P(E_{x'}^j;\RR)=|D\widetilde{\mathbf 1_{E_{x'}^j}}|(\RR)
\ge \frac{2}{L_j}\abs{E_{x'}^j}
\]
for almost every slice, including the null-slice and full-measure-slice cases.
Equivalently, the slicing identity can be written as
\[
\int_{R'_j}|D\widetilde{\mathbf 1_{E_{x'}^j}}|(\RR)\,\dd x'
=|D_{x_j}\mathbf 1_E|(\RR^k)
=\int_{\partial E}\abs{(\nu_E)_j}\,\dd\Haus^{k-1}.
\]
Applying the one-dimensional estimate and Fubini yields
\[
\int_{\partial E}\abs{(\nu_E)_j}\,\dd\Haus^{k-1}
\ge
\frac{2}{L_j}\Vol(E).
\]
Summing over $j$ gives \eqref{eq:box-anisotropic-ineq}. Equality holds for
$E=R$, since
\[
P_1(R)=2\Vol(R)\sum_{j=1}^k L_j^{-1}.
\]
This proves the asserted value of the supremum.

The displayed affine field satisfies
\[
\Div V_R=
2m_R\sum_{j=1}^k L_j^{-1}=1,
\]
so it is admissible. Its $j$th component ranges between $-m_R$ and $m_R$ on
$R$, whence
$\norm{V_R}_{L^\infty(R;\ell^\infty)}=m_R$. Thus $C_\infty(R)\le m_R$.
Lemma~\ref{lem:boundary-calibration} and the already computed supremum give
$C_\infty(R)\ge m_R$, so $C_\infty(R)=m_R$.
\end{proof}

\begin{theorem}[Quantitative boundary readout on axis-aligned boxes]
\label{thm:box-boundary-stability}
Here and below in this section, $k$ is the active dimension of the rectangular
box and of the top form $\omega_R$; it is not an additional inactive ambient
parameter from the preceding homotopy estimates.
Let
\[
R:=\prod_{j=1}^k[0,L_j],
\qquad
\omega_R:=\dd x_1\wedge\cdots\wedge\dd x_k,
\qquad
m_R:=\frac{1}{2\sum_{j=1}^kL_j^{-1}},
\]
and define
\[
\eta_R:=
2m_R\sum_{j=1}^k(-1)^{j-1}
\frac{x_j-L_j/2}{L_j}\,
\dd x_1\wedge\cdots\wedge\widehat{\dd x_j}
\wedge\cdots\wedge\dd x_k .
\]
Then $\dd\eta_R=\omega_R$ and
$\norm{\eta_R}_{\coeff,\infty}=m_R$. Let
$\eta\in\Omega^{k-1}_{\mathrm{sm}}(R)$, in the convention of
Section~2, meaning the restriction to $R$ of a smooth form on a neighborhood of
the closed box, satisfy $\dd\eta=\omega_R$, and set
$M:=\norm{\eta}_{\coeff,\infty}$. For
\[
F_{j,\lambda}:=\{x\in R:\ x_j=\lambda\},
\qquad
\lambda\in\{0,L_j\},
\]
write
\[
\iota_{j,\lambda}^\ast\eta
=
f_{j,\lambda}\,
\dd x_1\wedge\cdots\wedge\widehat{\dd x_j}\wedge\cdots\wedge\dd x_k
\]
with respect to the coordinate orientation on the face, and set
\[
s_{j,L_j}:=(-1)^{j-1},
\qquad
s_{j,0}:=-(-1)^{j-1}.
\]
Then
\begin{equation}
\label{eq:box-boundary-deficit}
\sum_{j=1}^k\sum_{\lambda\in\{0,L_j\}}
\int_{F_{j,\lambda}}\bigl(M-s_{j,\lambda}f_{j,\lambda}\bigr)
\,\dd\Haus^{k-1}
=
P_1(R)(M-m_R).
\end{equation}
Consequently
\begin{equation}
\label{eq:box-boundary-trace-stability}
\sum_{j=1}^k\sum_{\lambda\in\{0,L_j\}}
\norm{\iota_{j,\lambda}^\ast(\eta-\eta_R)}_{L^1(F_{j,\lambda};\coeff)}
\le
2P_1(R)(M-m_R).
\end{equation}
In particular, $M\ge m_R$, and every minimizer satisfies
\[
\iota_{j,\lambda}^\ast\eta
=
s_{j,\lambda}m_R\,
\dd x_1\wedge\cdots\wedge\widehat{\dd x_j}\wedge\cdots\wedge\dd x_k
\]
on every face.
\end{theorem}

\begin{proof}
The formula for $\eta_R$ is the form corresponding to the affine vector field
$V_R$ in Lemma~\ref{lem:box-self-cheeger} under the identification
\[
\eta=\sum_{j=1}^k(-1)^{j-1}V_j\,
\dd x_1\wedge\cdots\wedge\widehat{\dd x_j}\wedge\cdots\wedge\dd x_k.
\]
Thus $\dd\eta_R=(\Div V_R)\omega_R=\omega_R$ and
$\norm{\eta_R}_{\coeff,\infty}=\norm{V_R}_{L^\infty(R;\ell^\infty)}=m_R$.

Write the given $\eta$ in the same way, with associated vector field $V$.
Then $\Div V=1$ on $R$ and $M=\norm{V}_{L^\infty(R;\ell^\infty)}$.
On $F_{j,L_j}$, the outward normal is $e_j$ and
$V\cdot\nu=V_j=s_{j,L_j}f_{j,L_j}$; on $F_{j,0}$, the outward normal is
$-e_j$ and again $V\cdot\nu=s_{j,0}f_{j,0}$. Since
$\abs{\nu}_1=1$ on every face, Proposition~\ref{prop:quantitative-boundary-saturation}
applied to $\Omega=R$ gives \eqref{eq:box-boundary-deficit}.

For the affine minimizer one has
\[
\iota_{j,\lambda}^\ast\eta_R
=s_{j,\lambda}m_R\,
\dd x_1\wedge\cdots\wedge\widehat{\dd x_j}\wedge\cdots\wedge\dd x_k.
\]
Pointwise on each face,
\[
\abs{f_{j,\lambda}-s_{j,\lambda}m_R}
=\abs{s_{j,\lambda}f_{j,\lambda}-m_R}
\le
\bigl(M-s_{j,\lambda}f_{j,\lambda}\bigr)+(M-m_R),
\]
because $s_{j,\lambda}f_{j,\lambda}\le\abs{f_{j,\lambda}}\le M$ and
$M\ge m_R$ by Lemma~\ref{lem:box-self-cheeger}. Integrating over all faces and
using \eqref{eq:box-boundary-deficit} yields
\eqref{eq:box-boundary-trace-stability}. If $M=m_R$, the right-hand side
vanishes, so each continuous face coefficient
$f_{j,\lambda}-s_{j,\lambda}m_R$ vanishes $\Haus^{k-1}$-almost everywhere on
its face. A continuous function on a face that vanishes almost everywhere
vanishes everywhere, and hence every minimizer has the stated face trace on
every point of every face.
This is the pointwise upgrade from the almost-everywhere equality supplied by
Proposition~\ref{prop:quantitative-boundary-saturation}.
\end{proof}

\begin{corollary}[Deficit form of box trace stability]
\label{cor:box-delta-trace-stability}
Let $R$, $\omega_R$, $\eta_R$, and $m_R$ be as in
Theorem~\ref{thm:box-boundary-stability}. Suppose that
$\eta\in\Omega^{k-1}_{\mathrm{sm}}(R)$ satisfies $\dd\eta=\omega_R$ and that,
for some $\delta\ge0$,
\[
\norm{\eta}_{\coeff,\infty}\le m_R+\delta .
\]
Then
\begin{equation}
\label{eq:box-delta-trace-stability}
\sum_{j=1}^k\sum_{\lambda\in\{0,L_j\}}
\norm{\iota_{j,\lambda}^\ast(\eta-\eta_R)}_{L^1(F_{j,\lambda};\coeff)}
\le
2P_1(R)\delta .
\end{equation}
If $\delta=0$, the canonical face trace of
Theorem~\ref{thm:box-boundary-stability} is recovered on every coordinate
face.
\end{corollary}

\begin{proof}
Set $M:=\norm{\eta}_{\coeff,\infty}$. Theorem~\ref{thm:box-boundary-stability}
gives $M\ge m_R$ and
\[
\sum_{j=1}^k\sum_{\lambda\in\{0,L_j\}}
\norm{\iota_{j,\lambda}^\ast(\eta-\eta_R)}_{L^1(F_{j,\lambda};\coeff)}
\le
2P_1(R)(M-m_R).
\]
The assumed bound $M\le m_R+\delta$ implies $M-m_R\le\delta$, which proves
\eqref{eq:box-delta-trace-stability}. The last assertion is the case
$\delta=0$.
\end{proof}

\begin{theorem}[Sharp trace-error skeleton and minimizer kernel]
\label{thm:box-minimizer-structure}
Let $R=\prod_{j=1}^k[0,L_j]$, let $\eta_R$ and $m_R$ be as in
Theorem~\ref{thm:box-boundary-stability}, and let
$\eta\in\Omega^{k-1}_{\mathrm{sm}}(R)$ satisfy $\dd\eta=\omega_R$. For a
primitive $\eta$, write the face coefficients as in
Theorem~\ref{thm:box-boundary-stability}, set
\[
M:=\norm{\eta}_{\coeff,\infty},
\qquad
\delta:=M-m_R,
\qquad
q_{j,\lambda}:=s_{j,\lambda}f_{j,\lambda}-m_R
\quad\text{on }F_{j,\lambda}.
\]
Then $\delta\ge0$ and the signed trace error has the exact saturation-deficit
factorization
\begin{equation}
\label{eq:trace-skeleton-factorization}
q_{j,\lambda}=\delta-h_{j,\lambda},
\qquad
h_{j,\lambda}:=M-s_{j,\lambda}f_{j,\lambda}\ge0,
\end{equation}
with
\begin{equation}
\label{eq:trace-skeleton-deficit-mass}
\sum_{j=1}^k\sum_{\lambda\in\{0,L_j\}}
\int_{F_{j,\lambda}}h_{j,\lambda}\,\dd\Haus^{k-1}=P_1(R)\delta .
\end{equation}
Equivalently, the boundary trace lies in the following necessary capped
mean-zero skeleton:
\begin{equation}
\label{eq:trace-skeleton-mean-zero}
\sum_{j=1}^k\sum_{\lambda\in\{0,L_j\}}
\int_{F_{j,\lambda}}q_{j,\lambda}\,\dd\Haus^{k-1}=0,
\end{equation}
\begin{equation}
\label{eq:trace-skeleton-cap}
q_{j,\lambda}\le \delta\quad\text{pointwise on }F_{j,\lambda},
\end{equation}
\begin{equation}
\label{eq:trace-skeleton-lower-cap}
q_{j,\lambda}\ge -2m_R-\delta\quad\text{pointwise on }F_{j,\lambda},
\end{equation}
and
\begin{equation}
\label{eq:trace-skeleton-variation}
\sum_{j=1}^k\sum_{\lambda\in\{0,L_j\}}
\int_{F_{j,\lambda}}\abs{q_{j,\lambda}}\,\dd\Haus^{k-1}
=2
\sum_{j=1}^k\sum_{\lambda\in\{0,L_j\}}
\int_{F_{j,\lambda}}(q_{j,\lambda})_+\,\dd\Haus^{k-1}
\le 2P_1(R)\delta .
\end{equation}
Moreover, the exact minimizer fiber is completely described as follows:
$\eta$ is a minimizer,
\[
\norm{\eta}_{\coeff,\infty}=m_R,
\]
if and only if there is a smooth $(k-1)$-form
$\alpha\in\Omega^{k-1}_{\mathrm{sm}}(R)$ such that
\[
\eta=\eta_R+\alpha,
\qquad
\dd\alpha=0,
\qquad
\iota_{j,\lambda}^\ast\alpha=0
\quad(1\le j\le k,\ \lambda\in\{0,L_j\}),
\]
and
\[
\norm{\eta_R+\alpha}_{\coeff,\infty}\le m_R .
\]
Equivalently, under the identification
\[
\eta=\sum_{j=1}^k(-1)^{j-1}V_j\,
\dd x_1\wedge\cdots\wedge\widehat{\dd x_j}\wedge\cdots\wedge\dd x_k,
\]
the minimizers are precisely the fields
\[
V=V_R+W,
\qquad
\Div W=0,
\qquad
W\cdot\nu=0\ \text{on every coordinate face},
\qquad
\norm{V_R+W}_{L^\infty(R;\ell^\infty)}\le m_R .
\]
Finally, the cap in \eqref{eq:trace-skeleton-cap} is sharp in the class of
primitives with prescribed deficit: for every $a\in\RR$ and every
$1\le j\le k$, the primitive corresponding to the affine field $V_R+a e_j$
has deficit $|a|$ and has $q_{j,L_j}=a$, $q_{j,0}=-a$ on the two $j$-faces.
Hence no replacement of the upper cap $q\le\delta$ by $q\le c\delta$ with a
universal $c<1$ is possible. The displayed conditions are therefore the
optimal necessary trace skeleton obtained from the calibration; arbitrary
prescribed capped mean-zero boundary functions are not asserted to be
realizable.
\end{theorem}

\begin{proof}
The nonnegativity $\delta\ge0$ is the lower bound in
Theorem~\ref{thm:box-boundary-stability}. The functions
$h_{j,\lambda}=M-s_{j,\lambda}f_{j,\lambda}$ are nonnegative because
$s_{j,\lambda}f_{j,\lambda}\le\abs{f_{j,\lambda}}\le M$. Moreover,
$q_{j,\lambda}=s_{j,\lambda}f_{j,\lambda}-m_R=\delta-h_{j,\lambda}$, which is
\eqref{eq:trace-skeleton-factorization}. The boundary-deficit identity
\eqref{eq:box-boundary-deficit} gives \eqref{eq:trace-skeleton-deficit-mass}.

By the divergence theorem, or equivalently by summing the signed face-trace
coefficients in the proof of Theorem~\ref{thm:box-boundary-stability},
\[
\sum_{j=1}^k\sum_{\lambda\in\{0,L_j\}}
\int_{F_{j,\lambda}}s_{j,\lambda}f_{j,\lambda}\,\dd\Haus^{k-1}
=\int_R\omega_R
=\Vol(R).
\]
Since $P_1(R)=2\sum_j\prod_{i\ne j}L_i$ and
$m_R=(2\sum_jL_j^{-1})^{-1}$, one has $m_RP_1(R)=\Vol(R)$. Subtracting this
canonical identity gives \eqref{eq:trace-skeleton-mean-zero}. The pointwise
coefficient bound gives
\[
s_{j,\lambda}f_{j,\lambda}\le \abs{f_{j,\lambda}}\le M=m_R+\delta,
\]
which is exactly \eqref{eq:trace-skeleton-cap}. The opposite pointwise bound
$s_{j,\lambda}f_{j,\lambda}\ge-\abs{f_{j,\lambda}}\ge-M$ gives
\eqref{eq:trace-skeleton-lower-cap}. The mean-zero identity is an identity on
the disjoint boundary measure space
\[
\bigsqcup_{j=1}^k\bigsqcup_{\lambda\in\{0,L_j\}}
\bigl(F_{j,\lambda},\Haus^{k-1}\bigr),
\]
with $q$ equal to $q_{j,\lambda}$ on the corresponding component. Hence the
integral of $q$ over this disjoint union is zero, so its positive and negative
parts have equal integrals; this is the equality in
\eqref{eq:trace-skeleton-variation}. The cap
\eqref{eq:trace-skeleton-cap} and the identity
$\sum_{j,\lambda}\Haus^{k-1}(F_{j,\lambda})=P_1(R)$ give
\[
\sum_{j,\lambda}\int_{F_{j,\lambda}}(q_{j,\lambda})_+\,\dd\Haus^{k-1}
\le P_1(R)\delta,
\]
which proves the variation bound.

Assume first that $\eta$ is a minimizer and set
$\alpha:=\eta-\eta_R$. Since $\dd\eta=\dd\eta_R=\omega_R$, one has
$\dd\alpha=0$. The final assertion of
Theorem~\ref{thm:box-boundary-stability}, applied with $M=m_R$, gives the same
canonical trace for $\eta$ and $\eta_R$ on every coordinate face; hence
$\iota_{j,\lambda}^\ast\alpha=0$ for all $j$ and $\lambda$. The box constraint
is exactly $\norm{\eta_R+\alpha}_{\coeff,\infty}=\norm{\eta}_{\coeff,\infty}=m_R$.

Conversely, suppose that $\eta=\eta_R+\alpha$ satisfies the displayed
conditions. Then $\dd\eta=\omega_R$ and
$\norm{\eta}_{\coeff,\infty}\le m_R$. The lower bound
$\norm{\eta}_{\coeff,\infty}\ge m_R$ follows from
Lemma~\ref{lem:box-self-cheeger}, or equivalently from
Theorem~\ref{thm:box-boundary-stability}. Thus
$\norm{\eta}_{\coeff,\infty}=m_R$, so $\eta$ is a minimizer.

For the vector-field formulation write $\alpha$ in the same basis with
components $W_j$. Then $\dd\alpha=(\Div W)\omega_R$, so closedness of $\alpha$
is equivalent to $\Div W=0$. On the face $F_{j,\lambda}$ all pullbacks of basis
forms containing $\dd x_j$ vanish, and the only surviving component of
$\alpha$ is $(-1)^{j-1}W_j\,\dd x_1\wedge\cdots\wedge\widehat{\dd x_j}
\wedge\cdots\wedge\dd x_k$. Hence
$\iota_{j,\lambda}^\ast\alpha=0$ is equivalent to $W_j=0$ on that face, i.e.
to $W\cdot\nu=0$ there. The coefficient norm of $\eta_R+\alpha$ is the
$L^\infty(\ell^\infty)$ norm of $V_R+W$, giving the stated equivalence.

For sharpness, fix $a\in\RR$ and $j$. Let $V=V_R+a e_j$, and let $\eta$ be the
corresponding $(k-1)$-form under the displayed vector-field identification.
Since $a e_j$ is constant, $\Div V=1$, and hence $\dd\eta=\omega_R$. The $j$th
component of $V$ ranges over $[a-m_R,a+m_R]$, while every other component has
absolute value at most $m_R$. Thus
$\norm{\eta}_{\coeff,\infty}=m_R+|a|$ and the deficit is $|a|$. On
$F_{j,L_j}$ the outward flux of the perturbation $a e_j$ is $a$, while on
$F_{j,0}$ it is $-a$, so $q_{j,L_j}=a$ and $q_{j,0}=-a$. Taking $a>0$ attains
the upper cap on $F_{j,L_j}$, and taking $a<0$ attains it on $F_{j,0}$.
\end{proof}

\begin{corollary}[Local interface stability for one cell]
\label{cor:local-cell-interface-stability}
Let $R=\prod_{j=1}^k[0,L_j]$, let $\eta_R$ and $m_R$ be as in
Theorem~\ref{thm:box-boundary-stability}, and let
$\eta\in\Omega^{k-1}_{\mathrm{sm}}(R)$ satisfy $\dd\eta=\omega_R$. Suppose that,
for some $\delta\ge0$,
\[
\norm{\eta}_{\coeff,\infty}\le m_R+\delta .
\]
For each coordinate face $F_{j,\lambda}$, write
\[
\iota_{j,\lambda}^\ast(\eta-\eta_R)
=g_{j,\lambda}\,
\dd x_1\wedge\cdots\wedge\widehat{\dd x_j}\wedge\cdots\wedge\dd x_k,
\]
and let $\psi_{j,\lambda}\in L^\infty(F_{j,\lambda})$ be arbitrary face weights.
Define the linear one-cell interface observable
\[
\mathcal I_\psi(\eta):=
\sum_{j=1}^k\sum_{\lambda\in\{0,L_j\}}
\int_{F_{j,\lambda}}\psi_{j,\lambda}g_{j,\lambda}
\,\dd\Haus^{k-1} .
\]
Then
\begin{equation}
\label{eq:local-cell-interface-stability}
\abs{\mathcal I_\psi(\eta)}
\le
2P_1(R)\delta\,
\max_{j,\lambda}\norm{\psi_{j,\lambda}}_{L^\infty(F_{j,\lambda})} .
\end{equation}
In particular, every bounded linear face-flux readout used by a one-cell
evolution or balance-law update has the canonical value of $\eta_R$ at
deficit zero and has error bounded linearly by the interior saturation gap.
\end{corollary}

\begin{proof}
By Corollary~\ref{cor:box-delta-trace-stability},
\[
\sum_{j=1}^k\sum_{\lambda\in\{0,L_j\}}
\int_{F_{j,\lambda}}\abs{g_{j,\lambda}}\,\dd\Haus^{k-1}
\le 2P_1(R)\delta .
\]
Therefore
\[
\begin{aligned}
\abs{\mathcal I_\psi(\eta)}
&\le
\sum_{j=1}^k\sum_{\lambda\in\{0,L_j\}}
\norm{\psi_{j,\lambda}}_{L^\infty(F_{j,\lambda})}
\int_{F_{j,\lambda}}\abs{g_{j,\lambda}}\,\dd\Haus^{k-1} \\
&\le
\max_{j,\lambda}\norm{\psi_{j,\lambda}}_{L^\infty(F_{j,\lambda})}
\sum_{j=1}^k\sum_{\lambda\in\{0,L_j\}}
\int_{F_{j,\lambda}}\abs{g_{j,\lambda}}\,\dd\Haus^{k-1} \\
&\le
2P_1(R)\delta\,
\max_{j,\lambda}\norm{\psi_{j,\lambda}}_{L^\infty(F_{j,\lambda})} .
\end{aligned}
\]
If $\delta=0$, the same corollary gives $g_{j,\lambda}=0$ on every face, so
all such interface observables coincide with those of the affine minimizer
$\eta_R$.
\end{proof}
